% ===========================================================================
% 02 - Architektur: Service-Landschaft und Choreographie
% ===========================================================================

\begin{frame}{Service-Landschaft}
  \begin{columns}[T,onlytextwidth]
    \begin{column}{0.5\textwidth}
      \begin{itemize}
        \item \textbf{6 Microservices}, je ein Spring-Boot-Prozess.
        \item \textbf{2 zustandsbehaftet}: Configuration, Analysis Management.
        \item \textbf{4 stateless}: die Analyse-Services.
        \item Kommunikation ausschließlich über REST.
      \end{itemize}

      \vspace{0.3cm}
      \begin{block}{Fachlicher Schnitt}
        \small Jeder Service kapselt eine Business Capability --
        kein technischer Layer-Schnitt (ADR-01).
      \end{block}
    \end{column}
    \begin{column}{0.46\textwidth}
      \begin{center}
      \begin{tikzpicture}[
        node distance=4mm and 8mm,
        every node/.style={font=\scriptsize},
        svc/.style={draw=darkblue,rounded corners,fill=lightblue,minimum width=34mm,minimum height=6mm,align=center},
        data/.style={draw,ellipse,fill=lightgray,minimum width=26mm,minimum height=5mm,align=center},
        ana/.style={draw=darkblue,rounded corners,fill=white,minimum width=34mm,minimum height=6mm,align=center}
      ]
      \node[svc] (cfg) {configuration-service};
      \node[svc,below=of cfg] (am) {analysis-management-service};
      \node[data,below=of cfg] (cfgdb) {Configuration DB};
      \node[data,below=of am] (amdb) {Analysis DB};
      \node[ana,below=of amdb] (fluid) {fluid-analysis};
      \node[ana,below=of fluid] (thermal) {thermal-analysis};
      \node[ana,below=of thermal] (elec) {electrical-analysis};
      \node[ana,below=of elec] (ems) {engine-management};
      \draw[-{Latex}] (cfg) -- (cfgdb);
      \draw[-{Latex}] (am) -- (amdb);
      \draw[-{Latex},thick] (am) -- (fluid);
      \draw[-{Latex},thick] (fluid) -- (thermal);
      \draw[-{Latex},thick] (thermal) -- (elec);
      \draw[-{Latex},thick] (elec) -- (ems);
      \end{tikzpicture}
      \end{center}
    \end{column}
  \end{columns}
\end{frame}

\begin{frame}{Choreographie statt Orchestration}
  \begin{columns}[T,onlytextwidth]
    \begin{column}{0.5\textwidth}
      \textbf{Ablauf}
      \begin{enumerate}
        \item Analysis Management startet \textbf{nur} den Anchor \code{FLUID}.
        \item Jeder erfolgreiche Service ruft selbst den nächsten auf.
        \item Ergebnisse werden im \code{AnalysisCommand} mitgereicht.
        \item Engine Management berücksichtigt Vorergebnisse.
      \end{enumerate}

      \vspace{0.2cm}
      \begin{block}{Kernaussage}
        \small Kein zentraler Orchestrator -- die Ablaufsteuerung liegt
        verteilt in den Services (ADR-03).
      \end{block}
    \end{column}
    \begin{column}{0.46\textwidth}
      \begin{center}
      \begin{tikzpicture}[
        node distance=5mm,
        every node/.style={font=\scriptsize},
        stepbox/.style={draw=darkblue,rounded corners,fill=lightblue,minimum width=20mm,minimum height=8mm,align=center}
      ]
      \node[stepbox] (am) {Analysis\\Management};
      \node[stepbox,right=of am] (f) {Fluid};
      \node[stepbox,right=of f] (t) {Thermal};
      \node[stepbox,below=of t] (e) {Electrical};
      \node[stepbox,left=of e] (m) {Engine\\Management};
      \draw[-{Latex},thick] (am) -- node[above]{start} (f);
      \draw[-{Latex},thick] (f) -- (t);
      \draw[-{Latex},thick] (t) -- (e);
      \draw[-{Latex},thick] (e) -- (m);
      \end{tikzpicture}
      \end{center}

      \vspace{0.2cm}
      \begin{center}
      \begin{tikzpicture}[
        every node/.style={font=\scriptsize},
        cb/.style={draw=warnorange,rounded corners,fill=warnorange!10,minimum width=26mm,minimum height=8mm,align=center}
      ]
      \node[cb] {Circuit Breaker\\Thermal nicht erreichbar};
      \end{tikzpicture}
      \end{center}
    \end{column}
  \end{columns}
\end{frame}
